\begin{lemma}[Cotlar's inequality for $K_{\gamma,j}$, localized to a ball]
Fix $n \geq 1$ and $j \in \{1,\ldots,n\}$, let $\lambda$ be Lebesgue measure on
$\mathbb{R}^n$, let $\gamma_n$ be the Gaussian measure of
\noderef{GaussMeasure}, let $K_{\gamma,j}$ be the kernel of
\noderef{GaussRieszKernel} and let $M_2$ be the local maximal function of
\noderef{LocalMaximalFunction}.  Let
\[
  T : L^1(\gamma_n;\mathbb{R}) \longrightarrow
      \mathcal{M}\bigl(\gamma_n;\mathbb{R}^n\bigr)
\]
be \emph{any} linear map into the $\gamma_n$-almost-everywhere equivalence
classes of measurable $\mathbb{R}^n$-valued functions such that
\begin{enumerate}
\item[(T1)] $T(\iota v) = R_\gamma v$ for every $v \in L^2(\gamma_n)$, with
$\iota$ the inclusion of \noderef{ToL1} and $R_\gamma$ the transform of
\noderef{RieszTransform}; and
\item[(T2)] $t\,\gamma_n\bigl(\{x : |Tf(x)|_{\mathbb{R}^n} > t\}\bigr)
\leq 2\,\|f\|_{L^1(\gamma_n)}$ for every $f \in L^1(\gamma_n;\mathbb{R})$ and
every real $t > 0$.
\end{enumerate}
For $f \in L^1(\gamma_n)$, written also for its canonical representative, put
\[
  T_{j,\varepsilon}f(x) = \int_{|x-y|>\varepsilon}K_{\gamma,j}(x,y)f(y)\,dy ,
  \qquad
  T_{j,*}f(x) = \sup_{\varepsilon>0}\bigl|T_{j,\varepsilon}f(x)\bigr|
  \;\in\; [0,+\infty] ,
\]
the $y$-integral against Lebesgue measure and the supremum taken in the
extended nonnegative reals.  Let $R \geq 1$.  Then there is a constant
$C \in (0,\infty)$, depending only on $n$, $j$ and $R$, such that for every
$f \in L^1(\gamma_n)$ and \emph{every} $x$ with $|x| \leq R$,
\[
  T_{j,*}f(x)
    \;\leq\; C\Bigl[\,
      M_2\bigl(|(Tf)_j|^{1/2}\,\mathbf 1_{\overline B(0,R+2)}\bigr)(x)^2
      + M_2\bigl(f\,\mathbf 1_{\overline B(0,R+2)}\bigr)(x)
      + \|f\|_{L^1(\gamma_n)}\,\Bigr] ,
\]
the right-hand side read in $[0,+\infty]$, where $(Tf)_j$ denotes the $j$th
coordinate of the canonical representative of $Tf$.

\emph{Every $x$, not almost every $x$.}  The conclusion holds at every point of
the closed ball, with no exceptional set.  This matters downstream: the
consumer integrates the right-hand side in $x$ against weak-type bounds for
$M_2$, and an almost-everywhere conclusion in $x$ would still suffice there,
but the stronger form costs nothing and is what the argument produces.  The
exceptional sets that do arise --- three of them --- are all in the auxiliary
variable $z$ over which a point is selected inside the proof, and they are
discarded there.  Note also that both maximal functions and the level sets
below are unchanged if the representative of $Tf$ or of $f$ is modified on a
null set, so the choice of the canonical representative is a matter of fixing
notation and not an extra hypothesis.

\emph{Why both sides are read in $[0,+\infty]$.}  $M_2$ is extended-real valued
by \noderef{LocalMaximalFunction}, and $T_{j,*}f(x)$ is a supremum over
$\varepsilon > 0$ of real numbers which is not known in advance to be finite:
the bound $|T_{j,\varepsilon}f(x)| \leq C(\varepsilon)\|f\|_{L^1(\gamma_n)}$
available from the far majorant degrades as $\varepsilon \downarrow 0$.  Taking
the supremum in $[0,+\infty]$ rather than as a real supremum is what keeps the
inequality from being satisfiable by a conventional value where the family is
unbounded; it is the same discipline that makes \noderef{LocalMaximalFunction}
extended-real valued.  Each individual $T_{j,\varepsilon}f(x)$ is a genuine
absolutely convergent integral, by clause (1)(b) of
\noderef{GaussRieszKernelFarMajorant}.

\emph{The constant depends on $R$, necessarily, and the statement is
coordinatewise.}  $C$ degrades in $R$ through the far-majorant and
near-diagonal constants and through $a_{R+2} = \pi^{-n/2}e^{-(R+2)^2}$, the
lower bound for the Gaussian density on $\overline B(0,R+2)$; no uniformity in
$n$, $j$ or $R$ is claimed.  The object bounded is the coordinatewise maximal
truncation $\sup_\varepsilon|T_{j,\varepsilon}f|$, not the reference paper's
$\sup_\varepsilon|T_\varepsilon f|_{\mathbb{R}^n}$, which is a different
quantity; nothing here asserts anything about the latter.

\emph{The hypothetical style.}  (T1) and (T2) are written out in full and are
not imported by reference; the lemma is a statement about every $T$ satisfying
them, exactly as \noderef{CoordinateEndpointBound} and
\noderef{GaussRieszOffDiagonalL1} are, and it is not vacuous because
\noderef{L1Extension} exhibits such a $T$.
\end{lemma}
\begin{proof}
Fix $R \geq 1$ and $j$.  Write $p = \pi^{-n/2}$,
$a = \pi^{-n/2}e^{-(R+2)^2} > 0$, and $\omega_n = \lambda(\overline B(0,1))$,
which lies in $(0,\infty)$ because the closed unit ball is compact (so of
finite measure, shared context) and contains a nonempty open set (so of
positive measure).

Fix once and for all the following three constants.
\begin{itemize}
\item $C_a > 0$: the constant supplied by clause (1) of
\noderef{GaussRieszKernelFarMajorant} for the pair $(R,1)$.
\item $C_b > 0$: the constant supplied by the same clause for the pair
$(R+1,\tfrac12)$.
\item $C_h > 0$: the constant supplied by \noderef{GaussRieszKernelNearBounds}
at $R$ (its hypothesis $R > 0$ holds since $R \geq 1$).
\end{itemize}
Set
\[
  c = \max\Bigl(\,1,\ (C_a+C_b)\pi^{n/2},\
      C_h\omega_n2^n + 8\cdot2^n p\,a^{-1},\
      16\,p^2a^{-2}\,\Bigr) \in [1,\infty) ,
\]
which depends only on $n$, $j$ and $R$, and let $C$ be $c$ read in
$(0,\infty)$.  We show the asserted inequality with this $C$.

Fix $f \in L^1(\gamma_n)$ and $x$ with $|x| \leq R$, and abbreviate
\[
  \Phi = M_2\bigl(f\mathbf 1_{\overline B(0,R+2)}\bigr)(x),
  \quad
  \Psi = M_2\bigl(|(Tf)_j|^{1/2}\mathbf 1_{\overline B(0,R+2)}\bigr)(x),
  \quad
  N = \|f\|_{L^1(\gamma_n)} = \int|f|\,d\gamma_n < \infty .
\]
Write $w = f\mathbf 1_{\overline B(0,R+2)}$ and
$u = |(Tf)_j|^{1/2}\mathbf 1_{\overline B(0,R+2)}$ for the two functions whose
maximal functions these are.

\emph{Step 0: the infinite cases, and the reduction to a real bound.}
If $\Phi = +\infty$ or $\Psi = +\infty$ then the bracket on the right is
$+\infty$ and, since $C > 0$, the right-hand side is $+\infty$ and there is
nothing to prove.  So assume $\Phi,\Psi < \infty$ and write $\varphi,\psi \geq 0$ for the
corresponding reals.  The embedding of $[0,\infty)$ into $[0,+\infty]$ is an
order isomorphism onto its image and respects sums, products and squares, so
$\Psi^2 + \Phi + N$ is the image of the real $\psi^2+\varphi+N$ and
$C\,[\Psi^2+\Phi+N]$ is the image of $c(\psi^2+\varphi+N)$; and the supremum
defining $T_{j,*}f(x)$ is the supremum in $[0,+\infty]$ of the images of the
reals $|T_{j,\varepsilon}f(x)|$.  It therefore suffices to prove
\begin{equation*}
  \bigl|T_{j,\varepsilon}f(x)\bigr| \leq c\,\bigl(\psi^2 + \varphi + N\bigr)
  \qquad\text{for every } \varepsilon > 0 , \tag{$\star$}
\end{equation*}
and then take the supremum over $\varepsilon$ on the left.

\emph{Step 1: two consequences of finiteness of $\Phi$ and $\Psi$.}
By \noderef{LocalMaximalFunction}, for every $\rho \in (0,2]$
\[
  \frac{1}{\lambda(\overline B(x,\rho))}\int_{\overline B(x,\rho)}|w|\,d\lambda
    \leq \Phi ,
  \qquad
  \frac{1}{\lambda(\overline B(x,\rho))}\int_{\overline B(x,\rho)}|u|\,d\lambda
    \leq \Psi ,
\]
the integrals being lower Lebesgue integrals in $[0,+\infty]$ and
$\lambda(\overline B(x,\rho)) \in (0,\infty)$ for the same two reasons as
$\omega_n$ above.  Multiplying by that finite positive number and using
$\Phi,\Psi < \infty$,
\begin{equation*}
  \int_{\overline B(x,\rho)}|w|\,d\lambda \leq \varphi\,\lambda(\overline B(x,\rho)) < \infty,
  \qquad
  \int_{\overline B(x,\rho)}|u|\,d\lambda \leq \psi\,\lambda(\overline B(x,\rho)) < \infty
  \tag{1}
\end{equation*}
for every $\rho \in (0,2]$; in particular $w$ and $u$ are Lebesgue integrable
on every such ball.  We also record the dilation identity
$\lambda(\overline B(x,r)) = \omega_n r^n$ for $r \geq 0$, which holds for every
$n \in \mathbb{N}_0$: the closed ball $\overline B(x,r)$ is the image of
$\overline B(0,1)$ under $v \mapsto x + rv$, Lebesgue measure is translation
invariant (shared context) and satisfies $\lambda(rA) = r^n\lambda(A)$.

\emph{Step 2: the Gaussian/Lebesgue conversions used below.}
By \noderef{GaussMeasure}, $\gamma_n$ is $\lambda$ weighted by
$\varrho(y) = \pi^{-n/2}e^{-|y|^2}$, so for every measurable $v \geq 0$
\begin{equation*}
  \int v\,d\gamma_n = \int v(y)\varrho(y)\,d\lambda(y) . \tag{2}
\end{equation*}
Two consequences.  First, taking $v = |f|$ and using $\varrho(y) = p\,e^{-|y|^2}$,
\begin{equation*}
  \int_{\mathbb{R}^n}e^{-|y|^2}|f(y)|\,d\lambda(y) = \pi^{n/2}N < \infty .
  \tag{3}
\end{equation*}
The integrability of $y \mapsto e^{-|y|^2}|f(y)|$ against $\lambda$ together with the
value in $(3)$ is \noderef{GaussLebesgueL1DensityTransfer}, stated there for an
arbitrary element of $L^1(\gamma_n)$ and applied here to $f$; that node is what the
formalization invokes at this point.
Second, $\varrho \geq a$ on $\overline B(0,R+2)$ and $\varrho \leq p$
everywhere, so for every measurable $v \geq 0$ vanishing off
$\overline B(0,R+2)$,
\begin{equation*}
  a\int v\,d\lambda \leq \int v\,d\gamma_n \leq p\int v\,d\lambda . \tag{4}
\end{equation*}
In particular, with $v = |w|$, $\int_{\overline B(0,R+2)}|f|\,d\lambda
 \leq a^{-1}N < \infty$, so $f$ is Lebesgue integrable on
$\overline B(0,R+2)$.

\emph{Step 3: the case $\varepsilon \geq 1$.}
Let $\varepsilon \geq 1$.  By clause (1)(b) of
\noderef{GaussRieszKernelFarMajorant} at the pair $(R,1)$ --- applicable since
$f$ is $\gamma_n$-integrable and $|x| \leq R$ --- the function
$y \mapsto K_{\gamma,j}(x,y)f(y)$ is Lebesgue integrable on
$\{y : |x-y| > 1\}$, hence on the subset $\{y : |x-y| > \varepsilon\}$.  By
clause (1)(a) of the same node at the same pair, $|K_{\gamma,j}(x,y)| \leq
C_ae^{-|y|^2}$ whenever $|x| \leq R$ and $|x-y| \geq 1$, which holds throughout
that subset.  Hence, by $(3)$,
\[
  |T_{j,\varepsilon}f(x)|
    \leq \int_{\{|x-y|>\varepsilon\}}|K_{\gamma,j}(x,y)|\,|f(y)|\,d\lambda(y)
    \leq C_a\int e^{-|y|^2}|f|\,d\lambda
    = C_a\pi^{n/2}N \leq c\,N ,
\]
which is $(\star)$ for such $\varepsilon$.  For the rest of the proof fix
$\varepsilon$ with $0 < \varepsilon \leq 1$.

\emph{Step 4: the splitting of $f$, and the two inner balls.}
Put
\[
  B = \overline B(x,\varepsilon), \qquad B' = B(x,\varepsilon/2),
\]
and let $h, g \in L^1(\gamma_n)$ be the classes of the functions
$f\mathbf 1_{B}$ and $f\mathbf 1_{\mathbb{R}^n\setminus B}$ respectively; both
are $\gamma_n$-integrable because they are dominated by $|f|$.  Their
canonical representatives agree $\gamma_n$-almost everywhere, hence
$\lambda$-almost everywhere by clause (4) of
\noderef{GaussLebesgueBallComparison}, with those two functions, so we may and
do use $f\mathbf 1_B$ and $f\mathbf 1_{\mathbb{R}^n\setminus B}$ as the
representatives.  Since $f = f\mathbf 1_B + f\mathbf 1_{\mathbb{R}^n\setminus B}$
pointwise, $f = h + g$ in $L^1(\gamma_n)$.

Note $B' \subseteq B \subseteq \overline B(0,R+2)$, because $|x| \leq R$ and
$\varepsilon \leq 1$; and $g$ vanishes identically on $B$, hence
$\gamma_n$-almost everywhere on the open ball $B(x,\varepsilon)$.

\emph{Step 5: the two integral representations at $x$ and on $B'$.}
Apply \noderef{GaussRieszOffDiagonalL1} to $T$ (with its two properties, which
are (T1) and (T2) and hold by assumption), to $g$, to $j$, and with
$x_0 = x$ and $\rho = \varepsilon > 0$; its vanishing hypothesis was verified at
the end of Step 4.  Its clause (1) gives that for \emph{every}
$z \in B(x,\varepsilon/2) = B'$ the function
$y \mapsto K_{\gamma,j}(z,y)g(y)$ is Lebesgue integrable on $\mathbb{R}^n$, and
its clause (2) gives a $\gamma_n$-null set $N_1$ off which
\begin{equation*}
  z \in B' \implies (Tg)_j(z) = \int_{\mathbb{R}^n}K_{\gamma,j}(z,y)g(y)\,dy .
  \tag{5}
\end{equation*}
Since $\varepsilon > 0$ we have $x \in B'$, so clause (1) applies at $z = x$ as
well.  Moreover $\mathbb{R}^n\setminus B = \{y : |x-y| > \varepsilon\}$ exactly, so
$K_{\gamma,j}(x,\cdot)g$ is the restriction of $K_{\gamma,j}(x,\cdot)f$ to that
set, extended by zero; since it is Lebesgue integrable on $\mathbb{R}^n$ by the
previous paragraph,
\begin{equation*}
  T_{j,\varepsilon}f(x) = \int_{\mathbb{R}^n}K_{\gamma,j}(x,y)g(y)\,dy , \tag{6}
\end{equation*}
both sides being absolutely convergent.

\emph{Step 6: the difference estimate, uniformly over $z \in B'$.}
We claim that for every $z \in B'$,
\begin{equation*}
  \Bigl|\int K_{\gamma,j}(x,y)g(y)\,dy - \int K_{\gamma,j}(z,y)g(y)\,dy\Bigr|
    \leq C_h\omega_n2^n\varphi + (C_a+C_b)\pi^{n/2}N . \tag{7}
\end{equation*}
Fix $z \in B'$, so $|x-z| < \varepsilon/2 \leq 1/2$ and
$|z| \leq |x| + 1/2 \leq R+1$.  Both integrals in $(7)$ are absolutely
convergent by Step 5, and $g$ vanishes off
$\{y : |x-y| > \varepsilon\} = A_1 \cup A_2$, where
\[
  A_1 = \{y : \varepsilon < |x-y| \leq 1\}, \qquad A_2 = \{y : |x-y| > 1\} ,
\]
disjoint.  So the left-hand side of $(7)$ is at most
\[
  \int_{A_1}\bigl|K_{\gamma,j}(x,y)-K_{\gamma,j}(z,y)\bigr|\,|f(y)|\,d\lambda(y)
  \;+\;\int_{A_2}|K_{\gamma,j}(x,y)||f|\,d\lambda
  \;+\;\int_{A_2}|K_{\gamma,j}(z,y)||f|\,d\lambda .
\]
Here we used that the difference of the two absolutely convergent integrals is
$\int(K_{\gamma,j}(x,y)-K_{\gamma,j}(z,y))g(y)\,dy$, that $g$ vanishes off
$A_1 \cup A_2$ and equals $f$ on it, and that on $A_2$ the integrand is bounded
in absolute value by the sum of the two separate absolute values.  All three
integrals are finite: on $A_1$ the integrand is
$|K_{\gamma,j}(x,y)g(y) - K_{\gamma,j}(z,y)g(y)|$, dominated by
$|K_{\gamma,j}(x,\cdot)g| + |K_{\gamma,j}(z,\cdot)g|$, and on $A_2$ the two
integrands are $|K_{\gamma,j}(x,\cdot)g|$ and $|K_{\gamma,j}(z,\cdot)g|$; both
of those functions are Lebesgue integrable on $\mathbb{R}^n$ by Step 5.

\emph{The far part $A_2$.}  For $y \in A_2$: $|x| \leq R$ and $|x-y| \geq 1$,
so clause (1)(a) of \noderef{GaussRieszKernelFarMajorant} at $(R,1)$ gives
$|K_{\gamma,j}(x,y)| \leq C_ae^{-|y|^2}$; and
$|z-y| \geq |x-y| - |x-z| > 1 - \tfrac12 = \tfrac12$ with $|z| \leq R+1$, so
clause (1)(a) at $(R+1,\tfrac12)$ gives
$|K_{\gamma,j}(z,y)| \leq C_be^{-|y|^2}$.  With $(3)$, the two $A_2$ terms are
together at most $(C_a+C_b)\pi^{n/2}N$.

\emph{The near part $A_1$, by the H\"ormander estimate.}  Let $y \in A_1$.
Then $|x| \leq R$, $0 < \varepsilon < |x-y| \leq 1$ and
$|x-z| < \varepsilon/2 < |x-y|/2$, which are exactly the hypotheses of clause
(2) of \noderef{GaussRieszKernelNearBounds} at $R$; it gives
\[
  \bigl|K_{\gamma,j}(x,y)-K_{\gamma,j}(z,y)\bigr|
    \leq \frac{C_h\,|x-z|}{|x-y|^{\,n+1}}
    \leq \frac{C_h\,\varepsilon/2}{|x-y|^{\,n+1}} .
\]
Decompose $A_1$ dyadically.  For $k \in \mathbb{N}_0$ let
$D_k = \{y : 2^k\varepsilon < |x-y| \leq 2^{k+1}\varepsilon\}$; these are
pairwise disjoint, and every $y \in A_1$ lies in exactly one of them, namely
the one with $k$ the largest integer $\geq 0$ with $2^k\varepsilon < |x-y|$
(such a $k$ exists because $|x-y| > \varepsilon = 2^0\varepsilon$ and
$2^k\varepsilon \to \infty$).  Let $\mathcal{K}$ be the set of $k \in
\mathbb{N}_0$ with $D_k \cap A_1 \neq \emptyset$; for $k \in \mathcal{K}$ we
have $2^k\varepsilon < 1$, hence
\[
  r_k := 2^{k+1}\varepsilon < 2 ,
\]
so that $r_k \in (0,2]$ and $(1)$ applies at radius $r_k$.  This is the only
place the radius cap of \noderef{LocalMaximalFunction} is used, and it is used
with no margin: the cap is $2$ and the H\"ormander estimate is available only
for $|x-y| \leq 1$, which is exactly what forces $r_k < 2$.

For $k \in \mathcal{K}$, on $D_k$ we have $|x-y|^{-n-1} \leq
(2^k\varepsilon)^{-n-1}$ and $D_k \subseteq \overline B(x,r_k)$.  Also
$\overline B(x,r_k) \subseteq \overline B(0,R+2)$, since $|x| \leq R$ and
$r_k \leq 2$; hence $f = w$ on $\overline B(x,r_k)$ and, by $(1)$ and the
dilation identity of Step 1,
\[
  \int_{D_k}|f|\,d\lambda \leq \int_{\overline B(x,r_k)}|w|\,d\lambda
    \leq \varphi\,\omega_n r_k^{\,n} .
\]
Therefore
\[
  \int_{D_k\cap A_1}\bigl|K_{\gamma,j}(x,y)-K_{\gamma,j}(z,y)\bigr|\,|f|\,d\lambda
    \leq C_h\frac{\varepsilon}{2}\,(2^k\varepsilon)^{-n-1}\varphi\,\omega_n
        (2^{k+1}\varepsilon)^n
    = C_h\,\omega_n\,2^{\,n-1}\,2^{-k}\varphi ,
\]
because
$(2^k\varepsilon)^{-n-1}(2^{k+1}\varepsilon)^n = 2^{\,n-k}\varepsilon^{-1}$ and
$(\varepsilon/2)\cdot2^{\,n-k}\varepsilon^{-1} = 2^{\,n-k-1}$.  The sets $D_k \cap A_1$, $k \in \mathcal{K}$, are pairwise disjoint,
measurable and cover $A_1$, so by countable additivity the integral over $A_1$
is the sum of the integrals over them.  Summing the displayed bound over
$k \in \mathcal{K}$ and enlarging to all $k \geq 0$ (the terms are
nonnegative), $\sum_{k\geq0}2^{-k} = 2$ gives
\[
  \int_{A_1}\bigl|K_{\gamma,j}(x,y)-K_{\gamma,j}(z,y)\bigr|\,|f|\,d\lambda
    \leq C_h\omega_n2^{\,n}\varphi .
\]
Adding the two parts gives $(7)$.

\emph{Step 7: selection of a good $z$.}
Put
\[
  \lambda_1 = 8\cdot2^np\,a^{-1}\varphi \geq 0,
  \qquad
  \lambda_2 = \bigl(4p\,a^{-1}\psi\bigr)^2 \geq 0 ,
\]
and
\[
  E_1 = \{z \in B' : |(Th)_j(z)| > \lambda_1\},
  \qquad
  E_2 = \{z \in B' : |(Tf)_j(z)| > \lambda_2\} .
\]
We claim $\gamma_n(E_i) \leq \gamma_n(B')/4$ for $i = 1,2$.

First record two measure facts.  By the dilation identity,
$\lambda(B) = \omega_n\varepsilon^n$; and $B' = B(x,\varepsilon/2)$ differs
from $\overline B(x,\varepsilon/2)$ by the sphere of radius $\varepsilon/2$,
which is Lebesgue-null in $\mathbb{R}^n$ for $n \geq 1$, so
$\lambda(B') = \omega_n(\varepsilon/2)^n = 2^{-n}\lambda(B)$.  Also $B'$ is
measurable and $B' \subseteq \overline B(0,R+2)$, so clause (2) of
\noderef{GaussLebesgueBallComparison}, applied with $R+2$ in place of its $R$,
gives
\begin{equation*}
  \gamma_n(B') \geq a\,\lambda(B') = a\,2^{-n}\lambda(B) > 0 . \tag{8}
\end{equation*}

\emph{The set $E_1$.}  Suppose first $\varphi > 0$, so $\lambda_1 > 0$.  By
\noderef{CoordinateEndpointBound}, applied to $T$ with its two properties, to
the element $h$, to the coordinate $j$ and to $t = \lambda_1$,
\[
  \gamma_n(E_1) \leq \gamma_n\bigl(\{z : |(Th)_j(z)| > \lambda_1\}\bigr)
    \leq \frac{2\|h\|_{L^1(\gamma_n)}}{\lambda_1} .
\]
Now $\|h\|_{L^1(\gamma_n)} = \int_B|f|\,d\gamma_n$, and $B \subseteq
\overline B(0,R+2)$ so $f = w$ on $B$; by the upper half of $(4)$ and then
$(1)$ at radius $\varepsilon \in (0,2]$,
\[
  \|h\|_{L^1(\gamma_n)} \leq p\int_B|w|\,d\lambda \leq p\,\varphi\,\lambda(B) .
\]
Hence
\[
  \gamma_n(E_1) \leq \frac{2p\varphi\lambda(B)}{8\cdot2^np\,a^{-1}\varphi}
    = \frac{a\,2^{-n}\lambda(B)}{4} \leq \frac{\gamma_n(B')}{4}
\]
by $(8)$.  Suppose instead $\varphi = 0$.  Then by $(1)$ at radius
$\varepsilon$, $\int_B|w|\,d\lambda = 0$, so $w = 0$ and hence $f = 0$
$\lambda$-almost everywhere on $B$; therefore $f\mathbf 1_B = 0$
$\lambda$-almost everywhere, hence $\gamma_n$-almost everywhere by clause (4)
of \noderef{GaussLebesgueBallComparison}, so $h = 0$ in $L^1(\gamma_n)$.  By
linearity $Th = T0 = 0$, so $(Th)_j = 0$ $\gamma_n$-almost everywhere and
$E_1 = \{z \in B' : |(Th)_j(z)| > 0\}$ is $\gamma_n$-null.  In both cases
$\gamma_n(E_1) \leq \gamma_n(B')/4$.

\emph{The set $E_2$.}  Since $B' \subseteq \overline B(0,R+2)$, on $B'$ we have
$u = |(Tf)_j|^{1/2}$.  Note that $u$ is measurable, the canonical
representative of an almost-everywhere class being measurable, so $E_1$ and
$E_2$ are measurable sets.  Suppose first $\psi > 0$, so $\lambda_2 > 0$ and
$\lambda_2^{1/2} = 4p\,a^{-1}\psi$.  For $z \in B'$,
$|(Tf)_j(z)| > \lambda_2$ holds exactly when $u(z) > \lambda_2^{1/2}$, the
square root being increasing on $[0,\infty)$.  By Chebyshev's inequality for
$\lambda$ applied to the nonnegative function $u$ on $B'$ --- that is,
$\lambda_2^{1/2}\lambda(\{z \in B' : u(z) > \lambda_2^{1/2}\})
 \leq \int_{B'}u\,d\lambda$, which is immediate from
$\lambda_2^{1/2}\mathbf 1_{\{u>\lambda_2^{1/2}\}} \leq u$ on $B'$ --- together
with $B' \subseteq \overline B(x,\varepsilon/2)$ and $(1)$ at radius
$\varepsilon/2 \in (0,2]$,
\[
  \lambda(E_2) \leq \frac{1}{\lambda_2^{1/2}}\int_{\overline B(x,\varepsilon/2)}u\,d\lambda
    \leq \frac{\psi\,\lambda(\overline B(x,\varepsilon/2))}{\lambda_2^{1/2}}
    = \frac{\psi\,2^{-n}\lambda(B)}{4p\,a^{-1}\psi}
    = \frac{a\,2^{-n}\lambda(B)}{4p} .
\]
By clause (1) of \noderef{GaussLebesgueBallComparison}, which holds for every
set, $\gamma_n(E_2) \leq p\,\lambda(E_2) \leq a2^{-n}\lambda(B)/4 \leq
\gamma_n(B')/4$ by $(8)$.  Suppose instead $\psi = 0$.  Then by $(1)$ at
radius $\varepsilon/2$, $\int_{\overline B(x,\varepsilon/2)}u\,d\lambda = 0$,
so $u = 0$ $\lambda$-almost everywhere on $B'$, i.e.\ $(Tf)_j = 0$
$\lambda$-almost everywhere on $B'$, hence $\gamma_n$-almost everywhere there
by clause (4) of \noderef{GaussLebesgueBallComparison}; so
$E_2 = \{z \in B' : |(Tf)_j(z)| > 0\}$ is $\gamma_n$-null.  In both cases
$\gamma_n(E_2) \leq \gamma_n(B')/4$.

\emph{Two further null sets, both free.}  $N_1$ is the $\gamma_n$-null set of
$(5)$.  And since $T$ is linear, $Tg = Tf - Th$ as almost-everywhere classes,
so the canonical representative of $Tg$ agrees $\gamma_n$-almost everywhere
with the difference of those of $Tf$ and $Th$; let $N_2$ be the
$\gamma_n$-null set off which
\begin{equation*}
  (Tg)_j(z) = (Tf)_j(z) - (Th)_j(z) . \tag{9}
\end{equation*}

\emph{The good set is nonempty.}  Let
$G = B' \setminus (E_1 \cup E_2 \cup N_1 \cup N_2)$.  By subadditivity
(shared context),
\[
  \gamma_n(G) \geq \gamma_n(B') - \gamma_n(E_1) - \gamma_n(E_2)
    - \gamma_n(N_1) - \gamma_n(N_2)
    \geq \gamma_n(B') - \tfrac14\gamma_n(B') - \tfrac14\gamma_n(B')
    = \tfrac12\gamma_n(B') > 0
\]
by $(8)$, the subtraction being legitimate because
$\gamma_n(B') \leq p\,\lambda(B') < \infty$ by clause (1) of
\noderef{GaussLebesgueBallComparison}.
Hence $G \neq \emptyset$; fix $z \in G$.

\emph{Step 8: conclusion.}
By $(6)$, the triangle inequality, $(7)$ and then $(5)$ (legitimate since
$z \in B'$ and $z \notin N_1$),
\[
  |T_{j,\varepsilon}f(x)|
    \leq \Bigl|\int K_{\gamma,j}(x,y)g\,dy - \int K_{\gamma,j}(z,y)g\,dy\Bigr|
      + \Bigl|\int K_{\gamma,j}(z,y)g\,dy\Bigr|
    \leq C_h\omega_n2^n\varphi + (C_a+C_b)\pi^{n/2}N + \bigl|(Tg)_j(z)\bigr| .
\]
By $(9)$ (as $z \notin N_2$) and then $z \notin E_1 \cup E_2$,
\[
  |(Tg)_j(z)| \leq |(Tf)_j(z)| + |(Th)_j(z)| \leq \lambda_2 + \lambda_1
    = 16p^2a^{-2}\psi^2 + 8\cdot2^np\,a^{-1}\varphi .
\]
Combining, and using the definition of $c$,
\[
  |T_{j,\varepsilon}f(x)|
    \leq \bigl(C_h\omega_n2^n + 8\cdot2^np\,a^{-1}\bigr)\varphi
      + 16p^2a^{-2}\psi^2 + (C_a+C_b)\pi^{n/2}N
    \leq c\,(\varphi + \psi^2 + N) .
\]
This is $(\star)$ for $0 < \varepsilon \leq 1$; with Step 3 it holds for every
$\varepsilon > 0$.  Taking the supremum over $\varepsilon > 0$ in
$[0,+\infty]$ gives
$T_{j,*}f(x) \leq c(\psi^2 + \varphi + N) = C\,[\Psi^2 + \Phi + N]$, which is
the assertion.
%
% FORMALIZATION NOTE (cut points for ProofFormalization, recorded while the
% argument is fresh; see the reviewer's standing instruction that decomposition
% is free once the children's Lean closes, and must not be done in the stating
% phase).  Three natural children, in dependency order:
%   (C1) Steps 1-2 as a single arithmetic/measure lemma: for |x| <= R and
%        rho in (0,2], lambda(closedBall x rho) = omega_n rho^n,
%        closedBall x rho subset closedBall 0 (R+2), the averaged bound (1),
%        and the two-sided conversion (4).  Pure measure theory, no kernel.
%   (C2) Step 6 as "GaussRieszCotlarDifference": for |x| <= R, 0 < eps <= 1 and
%        z in B(x,eps/2), the bound (7).  This is the dyadic/Hoermander half and
%        is the longest single piece; it needs only the two kernel-estimate
%        nodes and (C1).  Do NOT make it a tablet node in the stating phase:
%        it would be exposed to Substantiveness Clause 2 against
%        GaussRieszKernelNearBounds.
%   (C3) Step 7 as "a point of B(x,eps/2) off three bad sets exists", i.e. the
%        selection, parameterized by the two thresholds.  Its content is
%        Chebyshev plus the measure bookkeeping (8); it is independent of the
%        kernel entirely.
% Step 8 then assembles (C2) and (C3) in five lines.
\end{proof}
